%!TeX program = xelatex
\documentclass{SYSUReport}
\usepackage{tabularx}
\usepackage{booktabs}
\usepackage{float}
\usepackage{amsmath}
\usepackage[backend=biber,style=numeric,sorting=none]{biblatex}
\addbibresource{reference.bib}
\setlength{\headheight}{15pt}
\lstset{breaklines=true,basicstyle=\ttfamily\small,columns=fullflexible}

\headl{}
\headc{}
\headr{信息安全技术期末考查}
\lessonTitle{信息安全技术期末考查}
\reportTitle{Watermark for LLM 三种算法对比实验}
\stuname{待填写}
\stuid{待填写}
\inst{待填写}
\major{待填写}
\date{\today}

\begin{document}
\cover
\thispagestyle{empty}
\clearpage
\pagenumbering{arabic}
\setcounter{page}{1}

\begin{abstract}
大语言模型生成文本的自动化能力正在快速提升，随之而来的学术诚信、虚假信息生成和合成数据污染等问题，使文本水印成为信息安全领域的重要研究方向。本文围绕三种代表性 LLM 水印方案展开复现与对比：KGW 通过伪随机绿名单和 logit 偏置在生成阶段嵌入统计信号；SWEET 针对代码等低熵场景，只在高熵 token 上嵌入和检测水印；EWD 则在检测阶段引入熵加权统计量，为高熵 token 分配更高权重。实验实现了从数据读取、模型生成、水印嵌入、水印检测、perplexity 评估到结果作图的完整链路。正式云端实验使用 OPT-1.3B 生成、OPT-6.7B 计算 PPL，在 C4 RealNewsLike 500 条新闻文本和 HumanEval 低熵代码任务上评估三种方法。结果表明，三种方法均能有效区分 human 与 watermarked 文本；SWEET 在保持高检测性能的同时带来更低 PPL 增量；EWD 适合作为 KGW 生成文本的增强检测器，但本身不是独立生成水印方案。
\end{abstract}

\section{背景和动机}
生成式大语言模型已经能够完成新闻写作、课程作业、代码生成和问答等任务。模型输出越接近人类文本，越难通过人工方式判断来源，这会带来三个典型风险。第一，恶意用户可以批量生成虚假新闻、评论或社交媒体内容，增加内容审核难度。第二，学生或程序员可能提交模型生成文本或代码，造成学术诚信与版权归属争议。第三，如果未来训练数据中混入大量机器生成内容，模型训练和评测本身也可能被污染。因此，能够在文本生成阶段主动嵌入、在事后自动检测的水印机制，是一类具有现实意义的信息安全技术。

与事后分类器不同，水印方法并不是在文本生成后再尝试识别写作风格，而是在生成过程中主动改变采样分布，使输出文本携带可统计检验的隐藏信号。这种思路的优势是检测结论可以用显式统计量解释，例如绿名单 token 的数量是否显著高于随机文本；不足之处是水印必须参与生成流程，且水印强度过高时可能损害文本质量。本文关注的 KGW、SWEET 和 EWD 都属于这一类围绕 token 采样分布设计的水印方法，因此实验重点放在“水印强度、检测准确性、文本质量损失”三者之间的权衡。

传统数字水印常见于图像、音频和视频，其核心是在人类难以察觉的载体空间中嵌入可检测信号。LLM 文本水印的困难在于文本是离散 token 序列，任何 token 修改都可能改变语义、语法或代码功能。KGW \cite{kirchenbauer2023watermark} 提出了一个简单有效的方向：不直接修改已生成文本，而是在模型采样时轻微提高部分 token 的概率，让最终文本在统计上包含更多“绿名单 token”。SWEET \cite{lee2024sweet} 观察到代码生成等任务的 token 熵较低，强行偏置低熵 token 会破坏功能，偏置太弱又难以检测，于是提出只处理高熵 token。EWD \cite{lu2024ewd} 则认为检测阶段也应充分利用熵信息，不能把所有 token 一视同仁。三者形成从“基础生成水印”到“选择性生成水印”再到“熵加权检测”的发展脉络。

\section{方法原理}
\subsection{KGW 水印}
KGW 的基本思想是：给定前文 token 和密钥，通过哈希函数确定一个伪随机种子，再把词表划分为绿名单 $G$ 和红名单 $R$。绿名单比例记为 $\gamma$。在生成第 $t$ 个 token 时，模型先输出原始 logit 向量，然后对所有绿名单 token 的 logit 加上偏置 $\delta$，再进行采样。若原始分布较高熵，绿名单 token 的采样概率会明显高于 $\gamma$；若文本不是由该水印生成，则 token 落入绿名单的概率近似为 $\gamma$。

这里的关键不是固定某些词必须出现，而是轻微改变候选 token 的相对概率。对高熵位置而言，模型本来就有多个合理候选，给绿名单 token 加上 $\delta$ 往往不会明显破坏语义，却会累积出可检测的统计偏差。对低熵位置而言，模型几乎只有一个合理 token，例如代码缩进、括号、固定关键字或常见 API 片段，此时 logit 偏置要么不起作用，要么会把输出推向不自然甚至错误的 token。因此 KGW 的简单性也带来一个问题：它对所有位置一视同仁，难以区分“适合嵌入水印的位置”和“最好不要干预的位置”。

检测时，检测器根据同一密钥和前文 token 复原每个位置的绿名单，统计待检测文本中绿名单 token 数 $|s|_G$。对于长度为 $T$ 的文本，KGW 使用一侧 z 检验：
\[
z = \frac{|s|_G - \gamma T}{\sqrt{T\gamma(1-\gamma)}}.
\]
当 $z$ 超过阈值时，判定文本含有水印。论文中常用 $z>4$ 作为较严格阈值，对应较低误报概率。

KGW 检测端的另一个优势是轻量：只要知道密钥、绿名单生成规则和待检测文本，就能复现每个位置的绿名单，不需要重新运行原始生成模型。这一点适合平台侧或第三方审核场景。但如果要分析 token 熵、PPL 或与 SWEET/EWD 做同一框架下的比较，实验中仍需要加载语言模型来获得 logits 或计算质量指标。

\subsection{SWEET 水印}
KGW 在自然语言新闻生成等高熵场景中效果较好，但在代码生成中会遇到两类问题：低熵 token 往往由语法或上下文强约束决定，强行提高绿名单概率可能破坏代码正确性；同时，低熵 token 很难被水印偏置改变，继续把它们计入检测统计量会稀释 z-score。SWEET 因此引入熵阈值 $\tau$。在生成阶段，只有当当前位置模型分布熵 $H_t>\tau$ 时才执行 KGW 式绿名单 logit 偏置；在检测阶段，也只统计熵高于阈值的 token。

换言之，SWEET 把 KGW 的“全位置加偏置”改成了“选择性加偏置”。如果某个位置的分布非常尖锐，说明模型对下一个 token 很确定，SWEET 会跳过该位置，避免为了水印而破坏代码结构或固定表达。如果某个位置分布较平坦，说明有多个候选都较自然，此时再施加绿名单偏置，既更容易产生绿名单 token，也更不容易造成明显质量下降。这个设计使 SWEET 同时改动了生成端和检测端：生成端少干预低熵 token，检测端也不让低熵 token 稀释统计量。

设高熵 token 数为 $N^h$，其中绿名单 token 数为 $N^h_G$，SWEET 的检测统计量为：
\[
z = \frac{N^h_G - \gamma N^h}{\sqrt{N^h\gamma(1-\gamma)}}.
\]
这相当于过滤掉难以被水印影响的低熵位置，使检测更集中于水印真正可能改变采样结果的区域。

熵阈值 $\tau$ 的选择决定了 SWEET 的取舍。阈值过低时，SWEET 接近 KGW，质量保护有限；阈值过高时，能保护更多低熵位置，但可计分 token 数减少，检测统计量可能不稳定。本文没有额外搜索阈值，而是保持统一配置，让 KGW、SWEET 和 EWD 在相同实验框架下比较，这样可以突出三种方法机制本身的差异。

\subsection{EWD 检测}
EWD 的核心贡献不在于提出新的生成器，而在于改造检测统计量。KGW 给所有 token 相同权重，SWEET 给高熵 token 权重 1、低熵 token 权重 0。EWD 认为 token 对检测结论的影响应与熵连续相关：高熵 token 更容易被水印偏置改变，因此应在检测中占更高权重；低熵 token 即使不在绿名单，也不应强烈削弱水印判断。

本文实现采用 EWD 论文中的线性权重形式。对第 $i$ 个被检测 token，根据模型分布计算 spike entropy $SE_i$，再令权重 $W_i=\max(SE_i-C_0,0)$。检测统计量为：
\[
z' = \frac{|s|_G - \gamma\sum_i W_i}{\sqrt{\gamma(1-\gamma)\sum_i W_i^2}},
\]
其中 $|s|_G$ 表示绿名单 token 权重和。EWD 可以用于 KGW 或 SWEET 生成的文本；本实验为了隔离变量，默认让 EWD 复用 KGW 生成文本，只替换检测器。

这种设置有助于避免把“生成质量变化”和“检测统计量变化”混在一起。由于 EWD 与 KGW 使用同一批生成文本，它们的 PPL 完全相同；如果 EWD 的 AUROC、TPR 或 z-score 更高，差异就来自检测端的熵加权，而不是来自生成文本本身更容易检测。EWD 的代价是检测时需要模型 logits 来计算 token 熵，因此部署成本高于只复现绿名单的 KGW 检测器。

\section{实验设计}
\subsection{模型和数据}
正式实验遵循统一模型原则。生成模型配置为 \texttt{facebook/opt-1.3b}，PPL/oracle 模型配置为 \texttt{facebook/opt-6.7b}，OPT 模型来自 Open Pre-trained Transformer 系列 \cite{zhang2022opt}。检测和生成均使用同一密钥、同一绿名单比例 $\gamma=0.25$、同一 logit 偏置 $\delta=2.0$。

选择统一模型而不是分别复现每篇论文的全部模型设置，是为了控制变量。KGW 原论文主要关注通用文本生成，SWEET 原论文使用代码模型并评估 pass@1，EWD 论文同时比较低熵代码和高熵文本。若完全照搬三篇论文的模型与数据，实验差异会同时来自算法、模型能力、数据分布和评测指标，难以判断某个指标变化究竟来自哪一部分。因此本文将三种算法放入同一实现框架中，对生成模型、检测密钥、采样参数和 PPL 计算方式进行统一，只比较水印生成器和检测器的差异。

主实验使用 C4 的 RealNewsLike 子集 \cite{raffel2020c4}，共 500 条。每条样本文本被切分为 prompt 和 human completion，prompt 输入生成模型，生成长度设为 200 tokens；同一条 prompt 分别生成无水印文本、KGW 水印文本和 SWEET 水印文本，EWD 对 KGW 水印文本进行加权检测。

补充实验改用 HumanEval 低熵代码任务。HumanEval 每条样本包含函数签名、自然语言说明、示例和参考实现，更接近 SWEET 论文中的真实代码生成评测。项目早期使用过 200 条合成结构化 Python 样本，但该数据集区分度过强，三种方法的检测指标容易同时饱和，因此不再作为低熵对比的主要数据集。

在数据组织上，每条样本都包含 \texttt{prompt} 和 \texttt{human\_completion}。human completion 作为负样本，用来估计误报；模型生成的水印 completion 作为正样本，用来估计检出率。HumanEval 的 prompt 通常已经包含函数签名、文档字符串和示例，这会让后续代码补全呈现更强的结构约束，因此比普通新闻文本更接近低熵场景。

\subsection{评价指标}
本文使用以下指标对比三种算法：
\begin{itemize}
    \item z-score：检测统计量，越高表示越可能含有水印。
    \item AUROC：将 human completion 作为负样本、水印生成文本作为正样本，衡量检测分数排序能力。
    \item TPR@5\%FPR：在误报率不超过 5\% 时的检测召回率。
    \item $z>4$ 检出率和误报率：对应 KGW 论文中常用的严格检测阈值。
    \item PPL：用 OPT-6.7B oracle 模型计算生成 completion 的困惑度，反映生成质量变化。
    \item 可计分 token 数和绿名单比例：用于解释水印嵌入强度和检测分数。
\end{itemize}

这些指标对应不同问题。AUROC 更关注排序能力，即水印文本是否整体排在人类文本前面；TPR@5\%FPR 更接近实际审核场景，因为误报过高会影响正常用户；$z>4$ 检出率则对应 KGW 论文中常用的固定阈值，便于解释统计显著性。PPL 不直接等于人类主观质量或代码正确率，但在统一 oracle 模型下可以反映生成分布是否被水印过度扭曲。对于 HumanEval，本文没有执行测试用例计算 pass@1，因此低熵实验的质量结论应理解为 PPL 层面的质量代理，而不是完整代码功能正确性结论。

\subsection{代码实现}
项目代码实现了完整实验链路。水印模块包含 KGW 和 SWEET 的生成处理器，可直接接入 HuggingFace 生成流程；同一模块还实现 KGW、SWEET 和 EWD 三类检测器。为保证检测和生成跨 CPU/GPU 一致，绿名单伪随机排列固定在 CPU 上生成，再搬运到目标设备。实验入口负责读取配置、加载模型、生成文本、计算检测分数和 PPL，并输出 CSV/JSONL；作图脚本根据结果绘制 z-score 分布、ROC 曲线和质量-可检测性权衡图。

实验输出分为三类文件：\texttt{generations.jsonl} 保存每个样本的 prompt、生成文本、PPL 和正样本 z-score，便于追查具体失败案例；\texttt{scores.csv} 保存 human 与 watermarked 两类文本的逐样本检测分数；\texttt{metrics.csv} 汇总 AUROC、TPR@5\%FPR、平均 z-score 和平均 PPL 等指标。EWD 在配置中通过 \texttt{reuse\_generation\_from: KGW} 复用 KGW 生成文本，因此同一批文本可以分别由 KGW 检测器和 EWD 检测器打分。

\section{实验结果与分析}
\subsection{C4 RealNewsLike 主实验}
表 \ref{tab:c4} 给出 C4 RealNewsLike 500 条主实验结果。三种方法都能有效区分 human 与 watermarked 文本，AUROC 均高于 0.995。在 $z>4$ 的固定阈值下，三种方法的误报率均为 0，检出率分别为 92.2\%、94.6\% 和 94.0\%。

C4 新闻文本整体熵较高，模型在许多位置存在多个自然候选，因此三种方法都能形成明显的绿名单偏置。KGW 的平均水印 z-score 为 8.839，已经远高于严格阈值；SWEET 在只选择高熵 token 的情况下仍达到 9.716，说明跳过低熵位置并没有削弱总体检测，反而减少了一些无效或有害的干预；EWD 的 9.669 则说明熵加权检测在高熵文本中也能保持较强排序能力。

\begin{table}[H]
\centering
\caption{C4 RealNewsLike 500 条主实验结果}
\label{tab:c4}
\small
\resizebox{\textwidth}{!}{
\begin{tabular}{lrrrrrrrr}
\toprule
算法 & 水印 z & human z & AUROC & TPR@5\%FPR & $z>4$ 检出 & $z>4$ 误报 & 平均 PPL & PPL 增量 \\
\midrule
KGW & 8.839 & -0.134 & 0.9955 & 0.990 & 92.2\% & 0\% & 21.388 & +7.986 \\
SWEET & 9.716 & -0.057 & 0.9982 & 0.994 & 94.6\% & 0\% & 17.239 & +3.837 \\
EWD & 9.669 & -0.056 & 0.9985 & 0.988 & 94.0\% & 0\% & 21.388 & +7.986 \\
\bottomrule
\end{tabular}
}
\end{table}

\begin{figure}[H]
    \centering
    \includegraphics[width=0.95\textwidth]{z_scores.png}
    \caption{C4 主实验中 human 与 watermarked z-score 分布}
    \label{fig:zscore}
\end{figure}

图 \ref{fig:zscore} 显示，human completion 的 z-score 集中在 0 附近，而水印文本整体明显高于 $z=4$ 阈值。KGW 和 SWEET 均存在少量低 z-score 离群点，主要来自过短生成或可计分 token 较少的样本；这说明固定阈值检测仍受文本长度影响。SWEET 的中位数 z-score 更高，同时 PPL 增量更低，说明选择性跳过低熵位置在高熵新闻场景中也能降低质量损失。

从误报角度看，三种方法在 C4 实验中的 human 文本均没有超过 $z>4$ 阈值，说明在当前参数下统计检验较保守。需要注意的是，固定阈值并不能保证所有正样本都被检出：当生成文本过短、重复、提前结束，或者可计分 token 较少时，即使文本由水印模型生成，z-score 也可能低于阈值。因此报告同时使用 AUROC、TPR@5\%FPR 和固定阈值检出率，而不只看单个指标。

\begin{figure}[H]
    \centering
    \includegraphics[width=0.72\textwidth]{roc_curves.png}
    \caption{C4 主实验 ROC 曲线}
    \label{fig:roc}
\end{figure}

图 \ref{fig:roc} 中三条 ROC 曲线均贴近左上角，说明排序意义上的检测能力很强。EWD 的 AUROC 最高，但 EWD 复用 KGW 生成文本，因此不能把它理解为独立生成算法优于 KGW；更准确的表述是：在 KGW 生成文本上，EWD 熵加权检测能略微改善检测分数排序。

\begin{figure}[H]
    \centering
    \includegraphics[width=0.72\textwidth]{quality_tradeoff.png}
    \caption{C4 主实验 PPL 与检测 z-score 权衡}
    \label{fig:tradeoff}
\end{figure}

图 \ref{fig:tradeoff} 展示了质量与可检测性的权衡。SWEET 位于更低 PPL、更高 z-score 的区域，说明它在本实验中取得了更好的质量-检测折中。KGW 和 EWD 的 PPL 相同，是因为 EWD 直接复用 KGW 的生成文本。C4 500 中存在一个极端样本只生成了极短 completion，导致 KGW/EWD PPL 达到 1385.43，并拉高平均 PPL；剔除 PPL 大于 100 的极端值后，KGW 平均 PPL 约为 18.655，SWEET 约为 16.965，SWEET 仍然更优。

\subsection{低熵代码补充实验}
表 \ref{tab:code} 给出 HumanEval 低熵代码补充实验结果。相比旧合成结构化代码数据，HumanEval 的检测指标不再饱和，更能体现 KGW、SWEET 和 EWD 在真实代码任务中的差异。需要说明的是，本实验仍使用 OPT-1.3B 作为统一生成模型，并以 PPL 作为质量代理指标；它不是 SWEET 论文中使用 StarCoder 并计算 pass@1 的完整代码正确性评测。

HumanEval 的结果更接近低熵场景的预期：所有方法仍然有效，但检测难度明显高于 C4。KGW 的 $z>4$ 检出率从 C4 的 92.2\% 降至 77.4\%，说明低熵代码 prompt 下，水印偏置并不总能稳定累积到足够高的 z-score。SWEET 通过跳过低熵位置降低了平均 PPL 增量，并把固定阈值检出率提高到 82.9\%。EWD 不改变生成文本，却把检出率提高到 92.1\%，说明在低熵检测时，连续熵权重比简单等权统计更能体现水印信号。

\begin{table}[H]
\centering
\caption{HumanEval 164 条低熵代码补充实验结果}
\label{tab:code}
\small
\resizebox{\textwidth}{!}{
\begin{tabular}{lrrrrrrrr}
\toprule
算法 & 水印 z & human z & AUROC & TPR@5\%FPR & $z>4$ 检出 & $z>4$ 误报 & 平均 PPL & PPL 增量 \\
\midrule
KGW & 5.820 & -1.917 & 0.9803 & 0.951 & 77.4\% & 0\% & 12.353 & +4.128 \\
SWEET & 6.418 & -0.417 & 0.9795 & 0.957 & 82.9\% & 0\% & 11.248 & +3.023 \\
EWD & 7.395 & -0.483 & 0.9884 & 0.988 & 92.1\% & 0\% & 12.353 & +4.128 \\
\bottomrule
\end{tabular}
}
\end{table}

\begin{figure}[H]
    \centering
    \includegraphics[width=0.95\textwidth]{code_z_scores.png}
    \caption{HumanEval 低熵代码中 human 与 watermarked z-score 分布}
    \label{fig:code_zscore}
\end{figure}

图 \ref{fig:code_zscore} 显示，HumanEval 中 watermarked 文本整体高于 human 文本，但三种方法的分离程度已有明显差别。KGW 的水印文本平均 z-score 为 5.820，严格阈值 $z>4$ 下检出率为 77.4\%；SWEET 将平均可计分 token 数从 KGW 的 120.88 降至 61.02，使检出率提高到 82.9\%，同时 PPL 增量从 +4.128 降至 +3.023；EWD 复用 KGW 生成文本，但通过熵加权检测把平均 z-score 提高到 7.395，AUROC 提高到 0.9884，$z>4$ 检出率达到 92.1\%。这说明 HumanEval 相比旧合成数据更能体现三种方法的取舍：SWEET 主要改善质量损失，EWD 主要增强检测统计量。

同时也可以看到，SWEET 的 human z-score 均值从 KGW 的 -1.917 上升到 -0.417，这是因为 SWEET 只检测高熵 token，human 文本中可计分位置更少，分数更接近 0。EWD 的 human 均值为 -0.483，仍保持较低误报，并且水印文本均值最高。换言之，EWD 并不是简单抬高所有文本的 z-score，而是更有效地扩大了 watermarked 与 human 的间隔。

\subsection{方法优缺点对比}
KGW 的优点是结构简单、检测不需要访问模型参数，只要知道密钥和哈希规则即可复原绿名单，因此部署成本低、统计解释清晰。缺点是低熵文本中水印信号可能被固定语法或上下文约束削弱，若提高 $\delta$ 又可能损害文本质量。

SWEET 解决的是“哪些 token 值得水印化”的问题：它跳过低熵位置，通常能降低对代码正确性或固定表达的破坏，同时提高有效统计样本中的绿名单比例。当前结果中，SWEET 取得了最低 PPL 增量，说明它在质量保持方面有优势。

EWD 解决的是“检测时怎样使用 token”的问题：它不再二值化地保留或删除 token，而是根据熵连续加权，因此比 KGW 更重视高熵 token，也比 SWEET 更少依赖人工阈值。不过 EWD 仍需要模型 logits，且本身不是独立生成水印方案。

综合来看，三种方法适合的使用场景并不完全相同。如果系统需要最简单、最容易部署的水印方案，KGW 是清晰的基线；如果生成内容对局部 token 极其敏感，例如代码、公式或固定格式文本，SWEET 更适合用于降低质量损失；如果平台已经能够在检测端访问模型 logits 或代理模型，EWD 可以作为更强的检测统计量叠加在 KGW 生成文本上。本文的实验也支持这一判断：SWEET 的主要收益体现在 PPL，EWD 的主要收益体现在检测分数。

\subsection{实验局限性}
本文实验仍有若干局限。首先，为了统一框架，HumanEval 实验使用 OPT-1.3B 生成代码，而不是专门的代码模型 StarCoder；因此生成代码的可执行正确性不能直接代表 SWEET 原论文中的 pass@1 结果。其次，PPL 只能作为文本质量的代理指标，尤其在代码任务中，低 PPL 并不必然意味着代码能通过测试。第三，当前实验只覆盖了无攻击场景，没有进一步测试改写、截断、拼接、变量重命名、回译等水印去除攻击。最后，SWEET 的熵阈值和 KGW 的 $\gamma,\delta$ 都会影响质量与检测的平衡，本文固定参数是为了便于比较，但不是最优参数搜索。

这些限制不影响本文的主要对比结论，但说明实验结果应被理解为“统一框架下的复现式比较”，而不是三篇论文完整实验设置的完全再现。若继续扩展，可以优先加入 MBPP 或 HumanEvalPack，并使用代码专用模型计算 pass@1；也可以把检测流程拆成不同攻击强度下的鲁棒性评测，从而更全面地评价水印在真实部署中的可靠性。

\section{运行效果与复现}
主实验命令如下：
\begin{lstlisting}[language=bash]
python scripts/prepare_assets.py --config configs/opt_full.yaml --hf-mirror --skip-models --dataset-samples 500 --dataset-output data/c4_realnewslike_500.jsonl
python scripts/run_experiment.py --config configs/opt_cloud_c4_500_ppl.yaml --output-dir results/opt_cloud_c4_500
python scripts/plot_results.py --results-dir results/opt_cloud_c4_500
\end{lstlisting}

低熵代码补充实验命令如下：
\begin{lstlisting}[language=bash]
python scripts/make_humaneval_dataset.py --output data/humaneval_164.jsonl
python scripts/run_experiment.py --config configs/opt_cloud_code_low_entropy_ppl.yaml --output-dir results/opt_cloud_code_low_entropy
python scripts/plot_results.py --results-dir results/opt_cloud_code_low_entropy
\end{lstlisting}

当前实现会同时加载 OPT-1.3B 生成模型和 OPT-6.7B PPL/oracle 模型。若云端显存不足，\texttt{device\_map: auto} 会触发 CPU offload，运行速度会明显下降。后续可改造为两阶段流程：先只加载 OPT-1.3B 完成生成和检测，再释放显存并加载 OPT-6.7B 补算 PPL。

\section{总结}
本文完成了 KGW、SWEET、EWD 三种 LLM 水印算法的复现式对比。三者的关系可以概括为：KGW 提供基础绿名单偏置框架，SWEET 在生成和检测阶段引入熵阈值选择，EWD 在检测阶段引入连续熵权重。云端正式实验使用 OPT-1.3B/OPT-6.7B，在 C4 RealNewsLike 500 条主实验和 HumanEval 低熵代码补充实验上验证水印检测的有效性。

实验结果表明，在高熵新闻文本中，三种方法均能稳定区分 human 与 watermarked 文本，且 $z>4$ 阈值下没有误报；在 HumanEval 低熵代码任务中，指标不再饱和，EWD 取得最高 AUROC 和严格阈值检出率，SWEET 则取得最低 PPL 增量。EWD 在本项目中应表述为“KGW 生成 + EWD 熵加权检测”，不能与 KGW、SWEET 完全并列为独立生成算法。

从安全角度看，文本水印的价值不在于给出绝对不可绕过的证明，而在于提供低成本、可量化、可审计的来源证据。未来可以进一步扩展 HumanEval/MBPP 等真实代码数据集，并增加同义改写、回译、截断、拼接攻击等鲁棒性实验。

\section*{小组贡献}
姓名与贡献待填写。建议提交前补充每位组员在论文阅读、算法实现、实验运行、报告撰写和视频录制中的具体工作。

\printbibliography
\end{document}
